% answers_short.tex —— 2013 年·四川卷·文科 · 简短答案速查
% 由 tools/make_short_answers.py 自 answers.tex 的 \gksolhead/\gkanswer 提取，
% 并与 2013/2013_数学真题_四川文_答案要点.md 逐题比对后生成，勿手改（改源 MD 后重跑）。

\begin{gknotice}
\textbf{使用说明}：本节只给最终答案便于快速核对，完整推导见解析版。答案为编写者独立推导，\textbf{非官方标准答案}。
\end{gknotice}

\gkpart{选择题}
\begin{center}\begin{tabular}{ccccccccccc}
\toprule
题号 & 1 & 2 & 3 & 4 & 5 & 6 & 7 & 8 & 9 & 10 \\
\midrule
答案 & B & D & B & C & D & A & A & C & C & A \\
\bottomrule
\end{tabular}\end{center}

\gkpart{填空题}
\begin{itemize}[leftmargin=2.6em]
  \item[\textbf{11.}] $1$。
  \item[\textbf{12.}] $2$。
  \item[\textbf{13.}] $36$。
  \item[\textbf{14.}] $\sqrt3$。
  \item[\textbf{15.}] $(2,4)$。
\end{itemize}

\gkpart{解答题}
\begin{itemize}[leftmargin=2.6em]
  \item[\textbf{16.}] 首项 $a_1=1$，公比 $q=3$，前 $n$ 项和 $S_n=\dfrac{3^n-1}2$。
  \item[\textbf{17.}] (1) $\dfrac45$；(2) $\dfrac{\sqrt2}{2}$
  \item[\textbf{18.}] (1) $P_1=\dfrac12$，$P_2=\dfrac13$，$P_3=\dfrac16$；(2) 甲为 $\dfrac{1027}{2100}$，$\dfrac{376}{2100}$，$\dfrac{697}{2100}$；乙为 $\dfrac{1051}{2100}$，$\dfrac{696}{2100}$，$\dfrac{353}{2100}$；乙所编写程序符合算法要求的可能性较大
  \item[\textbf{19.}] (1) 过 $P$ 作 $BC$ 的平行线即为 $l$；作图和证明见解析；(2) $\dfrac{\sqrt3}{6}$
  \item[\textbf{20.}] (1) $k<-\sqrt3$ 或 $k>\sqrt3$；(2) $n=\dfrac{\sqrt{15m^2+180}}5$（$-\sqrt3<m<\sqrt3$ 且 $m\neq0$）
  \item[\textbf{21.}] (1) 单调递减区间为 $(-\infty,-1)$，单调递增区间为 $(-1,0)$ 和 $(0,+\infty)$；(2) 证明见解析；(3) $a\in(-1-\ln2,+\infty)$
\end{itemize}
